\begin{proof}[Proof of \cref{obj:thm:log-free-bivariate-capacity}]
Fix integers \(n,d\geq2\), and set
\[
B:=\binom d2,\qquad K:=\lfloor n/4\rfloor.
\]
In particular, \(B>0\). Expectations of squares and sums of nonnegative terms are initially interpreted in the extended nonnegative reals.

\begin{enumerate}
\item
We use the expected absolute-supremum contraction inequality of
\citep[Theorem 12(4), with the complexity convention of Definition 2]{BartlettMendelson2002Complexities}
through \cref{obj:thm:full-design-lower,obj:lem:ht-transfer}. The outcome-law comparison uses the excess-variance identity of
\citep[Proposition 2.3, equation (2.4), under Assumption 2.1 and Definition 2.2]{Cytrynbaum2026Limits}
through \cref{obj:lem:published-prognostic-capacity}.

For every \(0<s\leq1\) and every admissible \(\pi\), \cref{obj:thm:full-design-lower} supplies
\[
c_s b_s(n,d)
\leq
\mathbb E_{m\sim\nu^\star_{n,d,s}}
\mathcal L_{n,d,s}(\pi,m),
\qquad
c_s b_s(n,d)\leq R_s(n,d).
\]
Also, directly from \cref{obj:def:criterion}, selecting any admissible kernel in the design infimum gives
\[
R_s(n,d)
\leq
\sup_{m\in\mathcal H_{d,s}}\mathcal L_{n,d,s}(\pi,m).
\]
We will apply this inequality to the kernel constructed below.
% lean: CausalSmith.Experimentation.BivariateSobolevDesignCapacity.full_design_lower
% lean: CausalSmith.Experimentation.BivariateSobolevDesignCapacity.published_prognostic_capacity
% lean: CausalSmith.Experimentation.BivariateSobolevDesignCapacity.log_free_bivariate_capacity

\item
We first construct the kernel using \(n,d\). For clarity, the spectral notation of \cref{obj:synth_3} is
\[
\mathbb T_2^d:=(\mathbb R/2\mathbb Z)^d,
\qquad
\operatorname{supp}(k):=\{j\in\{1,\ldots,d\}:k_j\ne0\},
\]
\[
t(k):=
\begin{cases}
\|k\|_2^2/|\operatorname{supp}(k)|,&k\ne0,\\
0,&k=0,
\end{cases}
\qquad
e_k(y):=\exp(\mathrm i\pi k^{\mathsf T}y)
\quad(k\in\mathbb Z^d,\ y\in\mathbb T_2^d).
\]
Let
\[
\mathscr K_d^+
:=
\left\{
k\in\mathbb Z^d:
1\leq|\operatorname{supp}(k)|\leq2,\
\text{the first nonzero coordinate of }k\text{ is positive}
\right\}.
\]
For \(h_{\mathrm{freq}}\in\mathbb N_0\), define \(k^{[h_{\mathrm{freq}}]}\) to be the
\((h_{\mathrm{freq}}+1)\)-st member of \(\mathscr K_d^+\), ordered by increasing \(t(k)\), with lexicographic ties. The prescribed real features are
\[
\phi_{2h_{\mathrm{freq}}+1}(y)
:=\sqrt2\,\operatorname{Re}e_{k^{[h_{\mathrm{freq}}]}}(y),
\qquad
\phi_{2h_{\mathrm{freq}}+2}(y)
:=\sqrt2\,\operatorname{Im}e_{k^{[h_{\mathrm{freq}}]}}(y).
\]
Since \(|e_k(y)|=1\),
\[
|\phi_h(y)|\leq\sqrt2
\qquad(h\geq1,\ y\in\mathbb T_2^d).
\]
For \(w=(w_1,\ldots,w_n)\in(\mathbb T_2^d)^n\), define
\[
a_h(w):=(\phi_h(w_i))_{i=1}^n,
\qquad
a(w):=(\phi_h(w_i))_{\substack{1\leq h\leq K\\1\leq i\leq n}}.
\]
Let
\[
Z^{\mathrm{sp}}(w,\zeta)
:=
\text{the output of ordered boundary rounding on }(a(w),\zeta),
\qquad
\zeta\in[0,1]^{3n},
\]
where the procedure is precisely \cref{obj:def:ordered-boundary-rounding}.

Draw independently a fair reflection array
\(\eta\in\{-1,1\}^{n\times d}\), a normalized Haar shift
\(T\in\mathbb T_2^d\), and a seed vector \(\zeta\) with independent uniform entries.
For every deterministic original sample
\(\mathbf x=(x_1,\ldots,x_n)\in([0,1]^d)^n\), define, modulo two,
\[
[\operatorname{lift}(\mathbf x,\eta)_i]_j
:=
\begin{cases}
(x_i)_j,&\eta_{ij}=1,\\
2-(x_i)_j,&\eta_{ij}=-1,
\end{cases}
\qquad
W_i(\mathbf x,\eta,T)
:=
\operatorname{lift}(\mathbf x,\eta)_i+T.
\]
The construction in \cref{obj:def:capacity-handle} is therefore
\[
\pi^\star_{n,d}(\mathbf x)
:=
\begin{cases}
\operatorname{Unif}(\{-1,1\}^n),&n\leq B,\\
\operatorname{Law}_{\eta,T,\zeta}
\bigl(Z^{\mathrm{sp}}(W(\mathbf x,\eta,T),\zeta)\bigr),&B<n.
\end{cases}
\]

For every fixed \(\mathbf x,\eta,T\), \cref{obj:lem:ordered-prefix-rounding}, applied with row bound \(\sqrt2\), gives
\[
\mathbb E_\zeta
Z_i^{\mathrm{sp}}(W(\mathbf x,\eta,T),\zeta)=0.
\]
Integration over \(\eta,T\), together with the independent-sign branch, yields
\[
\int_{\{-1,1\}^n}z_i\,\pi^\star_{n,d}(\mathbf x,dz)=0
\qquad(1\leq i\leq n)
\]
for every original sample. The feature evaluations, reflection lift, torus addition, and rounding output are Borel measurable; integrating their output against the probability seed law gives a Borel probability kernel. Drawing assignments from this covariate kernel gives
\[
Z\perp Y\mid X
\]
for every pre-assignment law allowed in \cref{obj:ass:assignment-independence}. Consequently,
\[
\pi^\star_{n,d}\in\mathcal D_{n,d,s}.
\]
Its branch choice, feature ordering, and seed distribution are independent of \(s\).
% lean: CausalSmith.Experimentation.BivariateSobolevDesignCapacity.piStar_pointwise_fair
% lean: CausalSmith.Experimentation.BivariateSobolevDesignCapacity.piStar_designClass

\item
We prove the upper guarantee for an arbitrary \(0<s\leq1\). First suppose \(n\leq B\). Conditional on \(X\), independent fair signs cancel every off-diagonal term:
\[
\mathbb E_Z\left[
\left(\sum_{i=1}^n Z_i m(X_i)\right)^2\Bigm|X
\right]
=
\sum_{i=1}^n m(X_i)^2.
\]
Thus, for every \(m\in\mathcal H_{d,s}\),
\[
\mathcal L_{n,d,s}(\pi^\star_{n,d},m)
=
4\int_{[0,1]^d}m(x)^2\,dx
\leq4,
\]
where the last inequality is supplied by \cref{obj:lem:exact-reflection-budget}.
Moreover,
\[
B/n\geq1
\quad\Longrightarrow\quad
(B/n)^s\geq1
\quad\Longrightarrow\quad
a_s(n,d)=1.
\]
Hence
\[
\mathcal L_{n,d,s}(\pi^\star_{n,d},m)
\leq4\leq3072a_s(n,d).
\]
% lean: CausalSmith.Experimentation.BivariateSobolevDesignCapacity.fairSignKernel_loss_le_four

\item
Suppose now that \(B<n\). Fix \(w\in(\mathbb T_2^d)^n\).
The moment guarantee in \cref{obj:lem:ordered-prefix-rounding} gives, for every \(h\geq1\),
\[
\mathbb E_\zeta
\left[
\bigl(a_h(w)^{\mathsf T}Z^{\mathrm{sp}}(w,\zeta)\bigr)^2
\right]
\leq32(\sqrt2)^2\min\{h,n\}
=
64\min\{h,n\}.
\]
This includes rows beyond the supplied prefix. When \(K=0\), the input matrix has zero rows and the same guarantee applies.

The two real rows associated with \(k^{[h_{\mathrm{freq}}]}\) satisfy the pointwise identity
\[
\left|
\sum_{i=1}^n
Z_i^{\mathrm{sp}}(w,\zeta)
e_{k^{[h_{\mathrm{freq}}]}}(w_i)
\right|^2
=
\frac{
\bigl(a_{2h_{\mathrm{freq}}+1}(w)^{\mathsf T}Z^{\mathrm{sp}}(w,\zeta)\bigr)^2
+
\bigl(a_{2h_{\mathrm{freq}}+2}(w)^{\mathsf T}Z^{\mathrm{sp}}(w,\zeta)\bigr)^2
}{2}.
\]
Indeed, the row imbalances are \(\sqrt2\) times the real and imaginary parts of the character imbalance. Taking expectations therefore gives
\[
\mathbb E_\zeta
\left|
\sum_{i=1}^n
Z_i^{\mathrm{sp}}(w,\zeta)
e_{k^{[h_{\mathrm{freq}}]}}(w_i)
\right|^2
\leq64\min\{2h_{\mathrm{freq}}+2,n\}.
\]
% lean: CausalSmith.Experimentation.BivariateSobolevDesignCapacity.character_imbalance_eq_real_rows
% lean: CausalSmith.Experimentation.BivariateSobolevDesignCapacity.spectral_rounding_character_second_moment

\item
We convert the row index into frequency complexity. Define, for \(t_{\mathrm{cut}}\geq1\),
\[
N_{\mathrm{real}}(t_{\mathrm{cut}})
:=
2\#\{k\in\mathscr K_d^+:t(k)\leq t_{\mathrm{cut}}\}.
\]
A support-one representative has a positive integer amplitude at most
\(\lfloor\sqrt{t_{\mathrm{cut}}}\rfloor\). A support-two representative has an ordered coordinate pair, a positive amplitude on its first coordinate, either sign on its second coordinate, and both absolute amplitudes at most
\(\lfloor\sqrt{2t_{\mathrm{cut}}}\rfloor\). Consequently,
\[
N_{\mathrm{real}}(t_{\mathrm{cut}})
\leq
2d\lfloor\sqrt{t_{\mathrm{cut}}}\rfloor
+
4B\lfloor\sqrt{2t_{\mathrm{cut}}}\rfloor^2.
\]
For \(t_{\mathrm{cut}}\geq1\),
\[
\lfloor\sqrt{t_{\mathrm{cut}}}\rfloor
\leq\sqrt{t_{\mathrm{cut}}}\leq t_{\mathrm{cut}},
\qquad
\lfloor\sqrt{2t_{\mathrm{cut}}}\rfloor^2
\leq2t_{\mathrm{cut}}.
\]
Also,
\[
2B=d(d-1)\geq d.
\]
It follows that
\[
N_{\mathrm{real}}(t_{\mathrm{cut}})
\leq2dt_{\mathrm{cut}}+8Bt_{\mathrm{cut}}
\leq12Bt_{\mathrm{cut}}.
\]

Every nonzero integer frequency satisfies
\[
\|k\|_2^2\geq|\operatorname{supp}(k)|,
\qquad t(k)\geq1.
\]
The first \(h_{\mathrm{freq}}+1\) representatives all have complexity at most
\(t(k^{[h_{\mathrm{freq}}]})\), so
\[
2h_{\mathrm{freq}}+2
\leq
N_{\mathrm{real}}\bigl(t(k^{[h_{\mathrm{freq}}]})\bigr)
\leq12B\,t(k^{[h_{\mathrm{freq}}]}).
\]
If \(Bt(k^{[h_{\mathrm{freq}}]})\leq n\), the preceding row-index bound gives a character moment at most \(768Bt(k^{[h_{\mathrm{freq}}]})\). If
\(n<Bt(k^{[h_{\mathrm{freq}}]})\), the bound by \(64n\) gives a character moment at most \(768n\). Thus
\[
\mathbb E_\zeta
\left|
\sum_{i=1}^n
Z_i^{\mathrm{sp}}(w,\zeta)
e_{k^{[h_{\mathrm{freq}}]}}(w_i)
\right|^2
\leq
768\min\{n,Bt(k^{[h_{\mathrm{freq}}]})\}.
\]

Every support-one or support-two frequency is a listed representative or its negative: its first nonzero coordinate determines the orientation, and there are finitely many representatives below each fixed complexity. Since real signs give
\[
\sum_i Z_i^{\mathrm{sp}}(w,\zeta)e_{-k}(w_i)
=
\overline{\sum_i Z_i^{\mathrm{sp}}(w,\zeta)e_k(w_i)},
\qquad
t(-k)=t(k),
\]
we obtain, for every supported frequency and every deterministic \(w\),
\[
\mathbb E_\zeta
\left|
\sum_{i=1}^n Z_i^{\mathrm{sp}}(w,\zeta)e_k(w_i)
\right|^2
\leq768\min\{n,Bt(k)\},
\qquad
1\leq|\operatorname{supp}(k)|\leq2.
\]
% lean: CausalSmith.Experimentation.BivariateSobolevDesignCapacity.featureFrequency_real_rows_le_complexity
% lean: CausalSmith.Experimentation.BivariateSobolevDesignCapacity.spectral_rounding_supported_second_moment

\item
Fix \(m\in\mathcal H_{d,s}\). Define its reflected function and Fourier coefficients by
\[
F(y):=m(b_d(y)),
\qquad
[b_d(y)]_j:=\min\{y_j,2-y_j\}
\quad(y\in[0,2)^d),
\]
\[
\widehat F(k)
:=
2^{-d}\int_{[0,2)^d}
F(y)\exp(-\mathrm i\pi k^{\mathsf T}y)\,dy
\qquad(k\in\mathbb Z^d).
\]
Under the independent uniform original sample, \cref{obj:lem:reflection-transfer} gives
\[
\operatorname{Law}(W,T)
=
\operatorname{Haar}(\mathbb T_2^d)^{\otimes n}
\otimes\operatorname{Haar}(\mathbb T_2^d)
\]
with normalized Haar laws. Its spectral identity gives
\[
\mathcal L_{n,d,s}(\pi^\star_{n,d},m)
=
\frac4n
\sum_{k\in\mathbb Z^d}
|\widehat F(k)|^2
\mathbb E_{W,\zeta}
\left|
\sum_{i=1}^n Z_i^{\mathrm{sp}}(W,\zeta)e_k(W_i)
\right|^2.
\]
Here \(W\) has the product Haar law, independently of the rounding seeds.

By \cref{obj:lem:exact-reflection-budget},
\[
\widehat F(0)=0,
\qquad
\widehat F(k)=0
\quad\text{if }|\operatorname{supp}(k)|>2,
\]
and
\[
\sum_{\substack{k\in\mathbb Z^d\\
1\leq|\operatorname{supp}(k)|\leq2}}
(1+\pi^2t(k))^s|\widehat F(k)|^2
\leq1.
\]
Integrating the conditional character bound and using
\[
\frac4n\,768\min\{n,Bt(k)\}
=
3072\min\{1,(B/n)t(k)\}
\]
therefore yields
\[
\mathcal L_{n,d,s}(\pi^\star_{n,d},m)
\leq
3072
\sum_{\substack{k\in\mathbb Z^d\\
1\leq|\operatorname{supp}(k)|\leq2}}
|\widehat F(k)|^2\min\{1,(B/n)t(k)\}.
\]
The integration and summation are justified by nonnegativity.
% lean: CausalSmith.Experimentation.BivariateSobolevDesignCapacity.spectralRounding_loss_fourier
% lean: CausalSmith.Experimentation.BivariateSobolevDesignCapacity.spectralRounding_loss_le_truncated

\item
We bound this truncated sum by the comparison scale. For every
\(x_{\mathrm{scalar}}\geq0\),
\[
\min\{1,x_{\mathrm{scalar}}\}\leq x_{\mathrm{scalar}}^s.
\]
Indeed, when \(0\leq x_{\mathrm{scalar}}\leq1\),
\[
x_{\mathrm{scalar}}
=x_{\mathrm{scalar}}^1
\leq x_{\mathrm{scalar}}^s,
\]
and when \(x_{\mathrm{scalar}}>1\),
\[
1\leq x_{\mathrm{scalar}}^s.
\]
These inequalities include \(s=1\).

Put
\[
r_{\mathrm{sample}}:=B/n.
\]
For every \(t_{\mathrm{weight}}\geq0\), the inequalities
\[
t_{\mathrm{weight}}\leq1+\pi^2t_{\mathrm{weight}},
\qquad
1\leq1+\pi^2t_{\mathrm{weight}}
\]
give
\[
\begin{aligned}
\min\{1,r_{\mathrm{sample}}t_{\mathrm{weight}}\}
&\leq
(r_{\mathrm{sample}}t_{\mathrm{weight}})^s\\
&=
r_{\mathrm{sample}}^s t_{\mathrm{weight}}^s\\
&\leq
r_{\mathrm{sample}}^s(1+\pi^2t_{\mathrm{weight}})^s,
\end{aligned}
\qquad
\min\{1,r_{\mathrm{sample}}t_{\mathrm{weight}}\}
\leq1\leq(1+\pi^2t_{\mathrm{weight}})^s.
\]
Use the former bound if \(r_{\mathrm{sample}}^s\leq1\), and the latter if
\(r_{\mathrm{sample}}^s>1\). This proves
\[
\min\{1,r_{\mathrm{sample}}t_{\mathrm{weight}}\}
\leq
\min\{1,r_{\mathrm{sample}}^s\}
(1+\pi^2t_{\mathrm{weight}})^s
=
a_s(n,d)(1+\pi^2t_{\mathrm{weight}})^s.
\]
Multiplying by the squared Fourier coefficients and summing gives
\[
\begin{aligned}
&\sum_{\substack{k\in\mathbb Z^d\\
1\leq|\operatorname{supp}(k)|\leq2}}
|\widehat F(k)|^2\min\{1,(B/n)t(k)\}\\
&\qquad\leq
a_s(n,d)
\sum_{\substack{k\in\mathbb Z^d\\
1\leq|\operatorname{supp}(k)|\leq2}}
(1+\pi^2t(k))^s|\widehat F(k)|^2
\leq a_s(n,d).
\end{aligned}
\]
Combining both design branches and then taking the fixed-function supremum proves
\[
\sup_{m\in\mathcal H_{d,s}}
\mathcal L_{n,d,s}(\pi^\star_{n,d},m)
\leq3072a_s(n,d).
\]
The construction was fixed before choosing \(s\), so this guarantee holds for the same kernel for every \(0<s\leq1\).
% lean: CausalSmith.Experimentation.BivariateSobolevDesignCapacity.spectral_min_rate_weight
% lean: CausalSmith.Experimentation.BivariateSobolevDesignCapacity.sobolev_truncated_spectral_sum_le_scale
% lean: CausalSmith.Experimentation.BivariateSobolevDesignCapacity.spectral_design_upper

\item
We also record the pathwise terminal-state guarantee. For
\(a\in\mathbb R^{K\times n}\), \(\zeta\in[0,1]^{3n}\), and
\(h_{\mathrm{move}}\in\mathbb N_0\), write
\[
u^{[h_{\mathrm{move}}]}(a,\zeta)
:=
\text{the fractional state after }h_{\mathrm{move}}
\text{ iterations of ordered boundary rounding},
\qquad
u^{[0]}(a,\zeta)=0.
\]
Applying \cref{obj:lem:ordered-prefix-rounding} to the rows
\(a_h(w)\), whose entries are bounded by \(\sqrt2\), gives
\[
u_i^{[n]}(a(w),\zeta)=Z_i^{\mathrm{sp}}(w,\zeta)
\qquad(1\leq i\leq n)
\]
for every \(w\) and every seed vector in the closed cube. In particular,
\[
u_i^{[n]}
\bigl(a(W(\mathbf x,\eta,T)),\zeta\bigr)
=
Z_i^{\mathrm{sp}}(W(\mathbf x,\eta,T),\zeta)
\]
for every original sample, every reflection array, every common shift, and every permitted rounding seed vector. Together with the displayed definition of \(\pi^\star_{n,d}\), this proves the two branch descriptions and the asserted equality between fractional coordinates and output signs.
% lean: CausalSmith.Experimentation.BivariateSobolevDesignCapacity.ordered_prefix_rounding
% lean: CausalSmith.Experimentation.BivariateSobolevDesignCapacity.log_free_bivariate_capacity

\item
Now fix \(0<s\leq1\). The two scales agree by their definitions:
\[
a_s(n,d)=b_s(n,d)=\min\{1,(B/n)^s\}\geq0.
\]
If \(n<B\), then \(B/n>1\), so
\[
a_s(n,d)=1.
\]
If \(B\leq n\), then \(0<B/n\leq1\), so
\[
a_s(n,d)=(B/n)^s.
\]
In particular, both expressions equal one at \(B=n\).

The constant calibration follows from
\[
0<48+32\pi^2,
\qquad
1\leq16384^s\leq16384,
\]
which imply
\[
c_1=\frac{2}{(48+32\pi^2)16384}>0,
\qquad
c_1\leq c_s.
\]
Furthermore,
\[
(48+32\pi^2)16384^s\geq48,
\qquad
c_s\leq\frac2{48}<3072.
\]
Multiplication by the nonnegative scale, the lower certificate, the design-infimum inequality, and the upper guarantee now give
\[
c_1a_s(n,d)
\leq c_sa_s(n,d)
\leq R_s(n,d)
\leq
\sup_{m\in\mathcal H_{d,s}}
\mathcal L_{n,d,s}(\pi^\star_{n,d},m)
\leq3072a_s(n,d).
\]
% lean: CausalSmith.Experimentation.BivariateSobolevDesignCapacity.cLower_uniform_positive
% lean: CausalSmith.Experimentation.BivariateSobolevDesignCapacity.log_free_bivariate_capacity

\item
For the prior in \cref{obj:def:cosine-prior}, its cutoff, feature dimension, and normalization are
\[
L:=\max\left\{1,\left\lceil\sqrt{4096n/B}\right\rceil\right\},
\qquad
M:=BL^2,
\qquad
C_0:=\sqrt{48+32\pi^2}.
\]
Since \(L\geq1\), \cref{obj:lem:cosine-prior} applies to every coefficient-sign array and gives
\[
\int_{[0,1]^d}m_\xi(x)\,dx=0,
\qquad
m_\xi\in\mathcal H_{d,s}.
\]
The coefficients are drawn independently of \(X\), as specified in \cref{obj:def:cosine-prior}. Using \(a_s=b_s\), the prior inequality from \cref{obj:thm:full-design-lower} becomes
\[
c_sa_s(n,d)
\leq
\mathbb E_{m\sim\nu^\star_{n,d,s}}
\mathcal L_{n,d,s}(\pi,m)
\qquad(\pi\in\mathcal D_{n,d,s}).
\]
% lean: CausalSmith.Experimentation.BivariateSobolevDesignCapacity.cosinePrior_legal_and_isotropic
% lean: CausalSmith.Experimentation.BivariateSobolevDesignCapacity.full_design_lower

\item
For the frozen uniform outcome class, \cref{obj:lem:published-prognostic-capacity} identifies the admissible kernels with \(\mathcal D_{n,d,s}\), identifies the prognostic image with \(\mathcal H_{d,s}\) modulo null sets, and gives
\[
\mathcal V_n(\pi,P)
=
\mathcal L_{n,d,s}(\pi,m_P)
\qquad
(\pi\in\mathcal D_{n,d,s},\ P\in\mathcal Q_{d,s}).
\]
Consequently, for every admissible \(\pi\),
\[
\sup_{P\in\mathcal Q_{d,s}}\mathcal V_n(\pi,P)
=
\sup_{m\in\mathcal H_{d,s}}\mathcal L_{n,d,s}(\pi,m).
\]
Taking the common design infimum proves
\[
\inf_{\pi\in\mathcal D_{n,d,s}}
\sup_{P\in\mathcal Q_{d,s}}\mathcal V_n(\pi,P)
=
R_s(n,d).
\]
Selecting \(\pi=\pi^\star_{n,d}\) in the kernel-level equality gives
\[
\sup_{P\in\mathcal Q_{d,s}}
\mathcal V_n(\pi^\star_{n,d},P)
=
\sup_{m\in\mathcal H_{d,s}}
\mathcal L_{n,d,s}(\pi^\star_{n,d},m).
\]
% lean: CausalSmith.Experimentation.BivariateSobolevDesignCapacity.published_prognostic_capacity

\item
Define the Gaussian subclass by
\[
\mathcal G_{d,s}
:=
\left\{
P:
P\text{ is the single-unit Gaussian completion of some }
m\in\mathcal H_{d,s}
\right\},
\]
using the completion in \cref{obj:synth_4}. For the completion of \(m\), its conditional means satisfy
\[
m_P=m\quad\text{almost everywhere},
\qquad
\mathbb E_P[O(1)-O(0)\mid U]
=
0.2+0.1\sqrt2\sin(2\pi U_1).
\]
Its two conditional potential-outcome variances are one. Since \(U_1\) is uniform and
\[
\int_0^1\sqrt2\sin(2\pi v)\,dv=0,
\qquad
\int_0^1\bigl(\sqrt2\sin(2\pi v)\bigr)^2\,dv=1,
\]
the law-specific benchmark is
\[
\mathsf V(P)=0.01+2\cdot1+2\cdot1=4.01.
\]

For every admissible \(\pi\) and every \(m\in\mathcal H_{d,s}\),
\cref{obj:lem:ht-transfer} supplies the actual-estimator identity and finiteness:
\[
n\operatorname{Var}(\widehat\theta)-4.01
=
\mathcal L_{n,d,s}(\pi,m)<\infty.
\]
Under this completion, \(\widehat\theta=\widehat T_P\). Taking the supremum over the completions of all \(m\in\mathcal H_{d,s}\) therefore gives
\[
\sup_{P\in\mathcal G_{d,s}}\mathcal V_n(\pi,P)
=
\sup_{m\in\mathcal H_{d,s}}\mathcal L_{n,d,s}(\pi,m).
\]
Taking the admissible-design infimum, and separately selecting the constructed kernel, yields
\[
\inf_{\pi\in\mathcal D_{n,d,s}}
\sup_{P\in\mathcal G_{d,s}}\mathcal V_n(\pi,P)
=R_s(n,d),
\]
\[
\sup_{P\in\mathcal G_{d,s}}
\mathcal V_n(\pi^\star_{n,d},P)
=
\sup_{m\in\mathcal H_{d,s}}
\mathcal L_{n,d,s}(\pi^\star_{n,d},m).
\]
Finally, \cref{obj:lem:published-prognostic-capacity} supplies the proper inclusion
\[
\mathcal G_{d,s}\subsetneq\mathcal Q_{d,s}.
\]
% lean: CausalSmith.Experimentation.BivariateSobolevDesignCapacity.gaussianWorst_eq_worstLoss
% lean: CausalSmith.Experimentation.BivariateSobolevDesignCapacity.gaussianCapacity_eq_capacity
% lean: CausalSmith.Experimentation.BivariateSobolevDesignCapacity.ht_transfer
% lean: CausalSmith.Experimentation.BivariateSobolevDesignCapacity.published_prognostic_capacity

\item
For every pair of real bounds
\[
0<\mathrm{lower}\leq1\leq\mathrm{upper},
\]
the bounded-density comparison in \cref{obj:lem:published-prognostic-capacity} gives
\[
R_s(n,d)
\leq
\inf_{\pi\in\mathcal D_{n,d,s}}
\sup_{P\in\mathcal A_{d,s}}\mathcal V_n(\pi,P),
\]
where \(\mathcal A_{d,s}\) has the density, moment, prognostic-budget, and permitted-intercept conditions specified there. This proves the finite-sample comparison for the full displayed range of density bounds.
% lean: CausalSmith.Experimentation.BivariateSobolevDesignCapacity.published_prognostic_capacity

\item
Let \(\varepsilon>0\), and suppose
\[
n\geq
B\max\left\{1,
\left(\frac{3072}{4.01\varepsilon}\right)^{1/s}\right\}.
\]
Since \(B,n,4.01\varepsilon\) are positive,
\[
\left(\frac{3072}{4.01\varepsilon}\right)^{1/s}
\leq\frac nB.
\]
Raising both sides to the positive power \(s\) gives
\[
\frac{3072}{4.01\varepsilon}
\leq\left(\frac nB\right)^s.
\]
The reciprocal-power identity is
\[
\left(\frac Bn\right)^s
\left(\frac nB\right)^s=1.
\]
Multiplying the preceding inequality by \((B/n)^s\) and then by
\(4.01\varepsilon\) yields
\[
3072a_s(n,d)
\leq3072(B/n)^s
\leq4.01\varepsilon.
\]
Hence the upper guarantee proves
\[
\sup_{m\in\mathcal H_{d,s}}
\mathcal L_{n,d,s}(\pi^\star_{n,d},m)
\leq4.01\varepsilon.
\]
% lean: CausalSmith.Experimentation.BivariateSobolevDesignCapacity.scale_sample_size_sufficient

\item
Conversely, suppose
\[
0<\varepsilon<c_s/4.01,
\qquad
R_s(n,d)\leq4.01\varepsilon.
\]
The lower bound gives
\[
c_sa_s(n,d)\leq4.01\varepsilon<c_s.
\]
If \((B/n)^s\geq1\), then \(a_s(n,d)=1\), contradicting this strict inequality. Thus
\[
(B/n)^s<1,
\qquad
a_s(n,d)=(B/n)^s.
\]
Multiplying the lower inequality by \((n/B)^s\), using the reciprocal-power identity, and dividing by \(4.01\varepsilon>0\) gives
\[
\frac{c_s}{4.01\varepsilon}
\leq\left(\frac nB\right)^s.
\]
Taking the positive \(1/s\)-power and multiplying by \(B>0\) proves
\[
n\geq
B\left(\frac{c_s}{4.01\varepsilon}\right)^{1/s}.
\]
% lean: CausalSmith.Experimentation.BivariateSobolevDesignCapacity.scale_sample_size_necessary

\item
For the asymptotic claims, fix \(0<s\leq1\) and a sequence \(d_n\geq2\). Work on the tail \(n\geq2\), and define
\[
B_n:=\binom{d_n}{2}.
\]
First,
\[
a_s(n,d_n)\longrightarrow0
\quad\Longleftrightarrow\quad
B_n/n\longrightarrow0.
\]
For the forward direction, a vanishing scale is eventually less than one, so eventually
\[
a_s(n,d_n)=(B_n/n)^s,
\qquad
B_n/n=a_s(n,d_n)^{1/s}.
\]
Continuity of the positive \(1/s\)-power at zero gives \(B_n/n\to0\).
For the reverse direction, continuity of the positive \(s\)-power gives
\[
(B_n/n)^s\longrightarrow0,
\qquad
0\leq a_s(n,d_n)\leq(B_n/n)^s,
\]
and squeezing gives the asserted scale convergence.

The pair count satisfies
\[
B_n=\frac{d_n(d_n-1)}2,
\qquad
\frac{d_n^2}{4}\leq B_n\leq\frac{d_n^2}{2},
\]
because \(d_n-1\geq d_n/2\) and \(d_n-1\leq d_n\). Therefore
\[
0\leq\frac{d_n^2}{n}\leq4\frac{B_n}{n},
\qquad
0\leq\frac{B_n}{n}\leq\frac{d_n^2}{n},
\]
and
\[
B_n/n\longrightarrow0
\quad\Longleftrightarrow\quad
d_n^2/n\longrightarrow0.
\]
Since \(d_n\geq0\) and \(n>0\) on this tail,
\[
\frac{d_n}{\sqrt n}
=
\sqrt{\frac{d_n^2}{n}},
\qquad
\left(\frac{d_n}{\sqrt n}\right)^2
=
\frac{d_n^2}{n}.
\]
Continuity of square roots and squares, together with the defining ratio criterion for little \(o\), gives
\[
B_n/n\longrightarrow0
\quad\Longleftrightarrow\quad
d_n=o(\sqrt n).
\]
% lean: CausalSmith.Experimentation.BivariateSobolevDesignCapacity.pairCount_quadratic_bounds
% lean: CausalSmith.Experimentation.BivariateSobolevDesignCapacity.scale_vanishing

\item
Suppose \(R_s(n,d_n)\to0\). For every \(n\geq2\), the lower certificate gives
\[
0\leq c_sa_s(n,d_n)\leq R_s(n,d_n).
\]
Squeezing and division by the fixed constant \(c_s>0\) yield
\[
c_sa_s(n,d_n)\longrightarrow0,
\qquad
a_s(n,d_n)\longrightarrow0.
\]
The scale equivalences just proved imply
\[
d_n=o(\sqrt n),
\qquad
B_n/n\longrightarrow0.
\]

Conversely, if \(B_n/n\to0\), then \(a_s(n,d_n)\to0\). The constructed admissible kernel gives, for every \(n\geq2\),
\[
0\leq R_s(n,d_n)
\leq
\sup_{m\in\mathcal H_{d_n,s}}
\mathcal L_{n,d_n,s}(\pi^\star_{n,d_n},m)
\leq3072a_s(n,d_n).
\]
The right-hand side tends to zero, so squeezing gives \(R_s(n,d_n)\to0\). Thus
\[
R_s(n,d_n)\longrightarrow0
\quad\Longleftrightarrow\quad
\frac{\binom{d_n}{2}}n\longrightarrow0
\quad\Longleftrightarrow\quad
d_n=o(\sqrt n).
\]
% lean: CausalSmith.Experimentation.BivariateSobolevDesignCapacity.dimension_necessity_of_scale_lower
% lean: CausalSmith.Experimentation.BivariateSobolevDesignCapacity.log_free_bivariate_capacity

\item
Finally, fix density bounds satisfying
\(0<\mathrm{lower}\leq1\leq\mathrm{upper}\).
For these fixed bounds and every sequence \(d_n\geq2\),
\cref{obj:lem:published-prognostic-capacity} supplies
\[
\inf_{\pi\in\mathcal D_{n,d_n,s}}
\sup_{P\in\mathcal A_{d_n,s}}\mathcal V_n(\pi,P)
\longrightarrow0
\quad\Longrightarrow\quad
d_n=o(\sqrt n).
\]
This establishes the bounded-density dimension implication under the stated conditions.
% lean: CausalSmith.Experimentation.BivariateSobolevDesignCapacity.published_prognostic_capacity
% lean: CausalSmith.Experimentation.BivariateSobolevDesignCapacity.log_free_bivariate_capacity
\end{enumerate}
\end{proof}
